\section{Conclusion}
We present a novel approach to modeling irregular time series data that can learn mechanisms for \emph{data-dependent missingness}. On all 3 large EHR datasets studied, our improved handling of missingness yields noticeably lower overall extrapolation error than several recent competitors.
While joint modeling of missingness indicators and observed features could be accomplished with many model architectures, by building upon \citet{schirmer2022modeling}'s  CRU backbone our implementation benefits from efficient computations and avoids the need for expensive ODE/SDE solvers.
Future work could investigate the even more flexible missing-not-at-random (MNAR) setting for irregular time series; 
though empirically every MNAR model has a MAR counterpart with equal fit 
\citep{molenbergsCounterpart2008} and thus the choice of MNAR requires belief in the underlying model assumptions not easily confirmed via empirical error analysis alone. More research is also necessary to ensure that our model's superior extrapolation capabilities translate effectively into enhanced clinical understanding and improved patient care. Additionally, while our comparisons are currently limited to other irregular time series models, examining how our approach stacks up against different types of models could provide broader insights. Lastly, identifying which specific vitals or lab results our methods most effectively predict remains an open question.